% ============================================================
% CONFIGURACIÓN GENERAL - APA 7 + adaptación institucional
% Motor: pdfLaTeX
% ============================================================

% ------------------------------------------------------------
% IDIOMA, CODIFICACIÓN Y FUENTE
% ------------------------------------------------------------
\usepackage[utf8]{inputenc}
\usepackage[T1]{fontenc}
\usepackage[spanish,es-nodecimaldot]{babel}

% Fuente de estilo Times compatible con pdfLaTeX.
% APA 7 admite Times New Roman 12 pt. Con pdfLaTeX, newtxtext/newtxmath
% es una alternativa Times muy cercana y portable.
\usepackage{newtxtext,newtxmath}

% ------------------------------------------------------------
% PÁGINA APA 7
% ------------------------------------------------------------
\usepackage[
  letterpaper,
  top=2.54cm,
  bottom=2.54cm,
  left=2.54cm,
  right=2.54cm,
  headheight=15pt,
  headsep=0.50cm
]{geometry}

\usepackage{setspace}
\doublespacing

\usepackage{ragged2e}
\AtBeginDocument{%
  \RaggedRight
  \setlength{\parindent}{1.27cm}% 1/2 pulgada
}
\setlength{\parskip}{0pt}
\hyphenpenalty=10000
\exhyphenpenalty=10000

% Si tu facultad exige texto justificado, comenta \RaggedRight arriba
% y usa \justifying.

% ------------------------------------------------------------
% NÚMERO DE PÁGINA ARRIBA A LA DERECHA
% ------------------------------------------------------------
\usepackage{fancyhdr}
\pagestyle{fancy}
\fancyhf{}
\fancyhead[R]{\thepage}
\renewcommand{\headrulewidth}{0pt}
\fancypagestyle{plain}{%
  \fancyhf{}
  \fancyhead[R]{\thepage}
  \renewcommand{\headrulewidth}{0pt}
}

% ------------------------------------------------------------
% TÍTULOS INSTITUCIONALES: CAPÍTULO 2, 2.1, 2.5.1, etc.
% ------------------------------------------------------------
% Nota: APA 7 puro no numera encabezados, pero el documento institucional
% de ejemplo sí los numera. Esta plantilla sigue ese ejemplo.
\usepackage{titlesec}

\titleformat{\chapter}[hang]
  {\normalfont\bfseries}
  {CAPÍTULO \thechapter.}
  {0.5em}
  {\MakeUppercase}
\titlespacing*{\chapter}{0pt}{0pt}{\baselineskip}

\titleformat{\section}[hang]
  {\normalfont\bfseries}
  {\thesection.}
  {0.5em}
  {\MakeUppercase}
\titlespacing*{\section}{0pt}{\baselineskip}{0pt}

\titleformat{\subsection}[hang]
  {\normalfont\bfseries}
  {\thesubsection.}
  {0.5em}
  {}
\titlespacing*{\subsection}{0pt}{\baselineskip}{0pt}

\titleformat{\subsubsection}[hang]
  {\normalfont\bfseries\itshape}
  {\thesubsubsection.}
  {0.5em}
  {}
\titlespacing*{\subsubsection}{0pt}{\baselineskip}{0pt}

\usepackage{indentfirst}

% ------------------------------------------------------------
% TABLAS Y FIGURAS APA
% ------------------------------------------------------------
\usepackage{graphicx}
\usepackage{booktabs}
\usepackage{array}
\usepackage{caption}
\usepackage{chngcntr}
\counterwithout{figure}{chapter}
\counterwithout{table}{chapter}

\addto\captionsspanish{%
  \renewcommand{\figurename}{Figura}%
  \renewcommand{\tablename}{Tabla}%
}

\captionsetup{
  justification=raggedright,
  singlelinecheck=false,
  labelsep=newline,
  labelfont=bf,
  textfont=it,
  font=normalsize,
  skip=6pt
}
\captionsetup[table]{position=above}
\captionsetup[figure]{position=above}

\newcommand{\notaAPA}[1]{%
  \par{\RaggedRight\noindent\textit{Nota.} #1\par}%
}

% ------------------------------------------------------------
% CITAS EN BLOQUE (40 PALABRAS O MÁS)
% ------------------------------------------------------------
\newenvironment{citaAPA}
{%
  \begin{list}{}{% 
    \setlength{\leftmargin}{1.27cm}%
    \setlength{\rightmargin}{0pt}%
    \setlength{\topsep}{0pt}%
    \setlength{\partopsep}{0pt}%
    \setlength{\parsep}{0pt}%
    \setlength{\itemsep}{0pt}%
  }
  \item[]\setlength{\parindent}{0pt}%
}
{\end{list}}

% ------------------------------------------------------------
% ENLACES Y URL
% ------------------------------------------------------------
\usepackage{xurl}
\usepackage[hidelinks]{hyperref}

% ------------------------------------------------------------
% REFERENCIAS APA 7
% ------------------------------------------------------------
\usepackage{csquotes}
\usepackage[
  backend=biber,
  style=apa,
  sorting=nyt,
  sortcites=true,
  uniquelist=false
]{biblatex}
\DeclareLanguageMapping{spanish}{spanish-apa}
\addbibresource{referencias.bib}

\setlength{\bibhang}{1.27cm}
\setlength{\bibitemsep}{0pt}

\defbibheading{apaheading}[REFERENCIAS]{%
  \clearpage
  \phantomsection
  \addcontentsline{toc}{chapter}{#1}%
  \begin{center}
    \textbf{#1}
  \end{center}
  \vspace{\baselineskip}
}

% ------------------------------------------------------------
% PORTADA APA PARA ESTUDIANTE
% ------------------------------------------------------------
\newcommand{\portadaAPA}[6]{%
  \thispagestyle{fancy}
  \vspace*{0.18\textheight}
  \begin{center}
    \textbf{#1}\par
    \vspace{2\baselineskip}
    #2\par
    #3\par
    #4\par
    #5\par
    #6\par
  \end{center}
  \vfill
  \clearpage
}
